\color{unimain}{\subsection{Устройство резервирования при отказе выключателя}}
\color{black}

\begin{enumerate}

\item
В устройстве предусмотрена функция резервирования при отказе выключателя (УРОВ). Данная функция предназначена для ликвидации повреждения в случае отказа выключателя поврежденного присоединения.

\item
Логика работы УРОВ выполнена с контролем протекания тока через выключатель. Значение по току срабатывания задается параметром \textbf{«Iср УРОВ В1(ЛВ)ВН»}. Функция УРОВ действует с регулируемой задержкой по времени \textbf{«Tср повт. УРОВ»} на повторное отключение своего выключателя, с задержкой \textbf{«Tср УРОВ»} на отключение смежных выключателей. 

\item
Предусмотрен подхват по току, включение которого производится программной накладкой \textbf{«УРОВ подхв.»}. 

\item
Действие УРОВ на смежные выключатели может быть выполнено с контролем сигнала РПВ. Ввод контроля сигнала РПВ в логику работы УРОВ осуществляется программной накладкой \textbf{«УРОВ контр. выкл.»}. 

\item 
Уставки УРОВ приведены в таблице \ref{app_set:urov}. Алгоритм работы УРОВ приведен на рисунке \ref{pict_urov:UROV}.

\medskip
\begin{figure}[H]
\centering
\includegraphics[width=1\textwidth,height=1\textheight,keepaspectratio]{\fbpath/09.01 БСК/761.3341 УРОВ/img/UROV.pdf}
\captionof{figure}{Функциональная схема логики УРОВ}\label{pict_urov:UROV}
\end{figure}


\end{enumerate} 